% Version 3.7: public receipt kernels, logical carriers and per-DATA acceptance.
\subsection{Published results and the global interpreter}
\label{sec:exports}
The Publish Layer carries complete regional results to Espresso.
The sender transports evidence; it does not choose the authoritative
global base or the grouping of atomic acceptance transactions.
An \tagg{EXPORT} retains its outer chain, version, expected-base and
interval fields for routing and diagnostics. For acceptance these are
\defterm{hints}. Signed result contents and authenticated source context
determine the chain, producing pair, ancestry and append permission.
A false hint neither authorizes invalid DATA nor vetoes a complete
valid result carried in the same record.

% Version 3.8: normalization preserves the original public result R0.
% Include within Published results, after its introduction and before extraction.
\paragraph{Logical EXPORT values and source denotations.}
\label{sec:logical-carrier}
The baseline EXPORT interface specifies a logical tuple
\[
 E=(c_E,\beta_E,b_E,I_E,W_E):
\]
chain identifier, branch version, expected global base, ordered DATA
interval and its evidence. The evidence contains the regional
receipts, ancestry, boundaries, historical authorities and native
context required by that interface. These are structured logical
values. No deployed binary encoder is fixed by this specification.
The signed objects retain their original bytes and hashes under
normalization; it does not reserialize signed content.

The logical structure provides field and list positions for DATA and
DRAIN bodies, signed proposals, certificates, ancestry, admission
witnesses, boundaries and source proofs. A complete result is the
DATA and direct or covering receipt with the required validation
objects specified by the original publication interface in
\cref{sec:delivery}, namely $R_0$ in \cref{sec:input-report-carriage}.
The new regional REPORT sidecar is not an extra field or required edge
of this old public value. Every original admission, proposal and source
object remains present; this is not a kernel-only carrier. Its occurrence
in these old tuple fields is determined before normalization is
defined. It can occur inside ancestry rather than in the sender's
declared output list. A wrong outer value or an invalid unrelated
proof does not change which objects those fields contain.

A reference retains its baseline denotation: the same content and
authenticated source identity that the original validation interface
uses. It is an edge to an object, not the object's bytes. Public source
objects can be represented inline or by aliases in the authenticated
retained dictionary. That dictionary retains authenticated raw
Espresso objects, including evidence from rejected records, and
authenticated L1 objects and proofs obtained through their own source
interface. Checkpoint objects are limited to the objects established
by that checkpoint's specified authentication contract.

An alias may identify an object by its content commitment or existing
authenticated source identity. An exact occurrence position may help
retrieve it, but a redundant explicit position is not a new validity
requirement. If several authenticated occurrences contain the same
object, derive a canonical representative from those occurrences.
An L1 authority or proof need not have been embedded in an earlier
Espresso record. The separate requirement that an Espresso selection
reference precede the containing EXPORT block remains unchanged.

Private caches implement this public object lookup; they do not add
new public denotations. A cache miss for an already specified object is
a retrieval task, while an absent required object is an incomplete
proof. A local regional cache reference can satisfy a point-to-point
validation bundle under its existing receiver-resolution condition;
that fact alone does not make a private-only object publicly supplied
by an EXPORT. This distinction follows the existing public
complete-result visibility obligation, rather than a new decoder's
success or failure.

\paragraph{Structural normalization.}
Write $\tagg{Flat}(E,\mathcal H)$ for the logical graph representation
of the tuple and its baseline authenticated source context $\mathcal H$.
It preserves every DATA and evidence field, including nested ancestry
and boundary evidence, in a fixed field/list order. Each typed object
keeps its original content, dependency edges and logical occurrence
path. Derive edges from the actual object fields, not a sender's
auxiliary index. Shared objects retain their incoming edges and
occurrence paths. References retain the same authenticated denotations
just described. This graph definition is not an instruction to fetch
every referenced target before processing the record.

The actual inventory is incremental. Traverse the supplied inline
structure and already retained authenticated objects, visiting each
actual object once. Preserve an unresolved reference as a typed edge;
do not wait for its target merely to finish this traversal. Extract
the required projection consisting of DATA bodies, FINAL signatures,
hash-linked paths and the fixed incoming context used by
$\tagg{CheckFINAL}$. Resolve references to these required objects through
the existing public input and retained-history interface. An optional
recovery authority, regional L1-reference proof or native sidecar is
not on this projection, so its acquisition cannot precede or gate
discovery of neighboring kernels. Inline structured children are still
inventoried even when their enclosing proof is invalid.

This distinction uses the baseline object meanings. An L1 authority
record denotes its specified authority fields, not a hidden regional
DATA/FINAL carrier; a native proof sidecar is not the missing DATA or
FINAL signatures. The complete-result transport interface supplies
the result's actual DATA and commitment evidence, with its required
ancestry, either in the logical record or through its existing public
source denotations. A kernel referenced in retained Espresso evidence
is read there. It is not obtained by first validating an optional
$\tagg{AUTH}_{r,g}$ or native claim. If a concrete compression hides
required DATA/FINAL bytes behind a new resolver, its compatibility and
access obligations must be proved separately; the logical graph lemma
does not manufacture them. No unresolved edge or cycle supplies a
missing required object.

Do not filter this traversal by outer chain, version, base or interval
hints, or stop it because a sibling or enclosing proof is invalid.
Each candidate still checks its own required objects. Preserve
alternative proofs of the same DATA separately so that one failed
proof does not erase another complete proof. The fixed logical
occurrence path, including a derived authenticated source occurrence
where needed, orders proof choices. A concrete framed encoding may
implement these paths but is not assumed by this definition.

\begin{proposition}[Preservation of logical result evidence]
\label{prop:carrier-preservation}
For a baseline logical EXPORT $E$ and its authenticated source context
$\mathcal H$, structural normalization preserves every complete
regional-result subgraph represented by its DATA and evidence fields,
with references interpreted under their unchanged baseline source
denotations. It therefore preserves that result's public FINAL kernel.
Changing only outer hints or adding unrelated invalid structured
evidence cannot remove the subgraph. Normalization creates no payload,
receipt, authority or append permission.
\end{proposition}
\begin{proof}
Fix the finite validation witness of a complete result in the original
logical fields. Induct over its object structure as a statement about
graph denotation. A leaf keeps its original content. For a compound
object every child remains represented independently of outer hints
and sibling validity; the induction hypothesis preserves the children
and their actual dependency edges. A shared object keeps its content
and additional occurrence edges. A reference denotes the same
authenticated source object before and after normalization, so the
corresponding case applies to that object. Retaining its unresolved
edge is sufficient for this representational statement; no acquisition
deadline for the target is inferred.
Thus every object and edge in the old witness survives. Its direct
receipt, DATA bodies and exact covering path in particular survive,
which preserves the public kernel of \cref{sec:receipt-kernel}.
For an actual complete public input, the DATA/FINAL projection supplied
by the inherited delivery interface is therefore the same projection
inventoried from its inline and retained source objects. Its extraction
does not require acquisition of the optional witness targets excluded
from $\tagg{CheckFINAL}$.
All signed content, authorities and original provenance are unchanged.
No traversal rule adds an absent leaf or uses an outer hint as an
inclusion condition. This proves both preservation and non-fabrication.
\end{proof}

This is a representation result. The revised interpreter deliberately
changes acceptance of bad-hint records; the proposition does not claim
that its acceptance trace equals the baseline's all-or-nothing EXPORT
trace. Finiteness of the inventory also supplies no processing-time
bound by itself.

\paragraph{Delivery scope and concrete encodings.}
The inherited Publish Layer and Espresso interfaces make the complete
logical result and its record available to the covered online,
history-following readers under their stated admission and load limits.
The original tuple fields and source denotations specify what is
delivered; \cref{prop:carrier-preservation} proves that normalization
preserves it. This does not define ``carrying'' to mean success of the
new extractor. A result already supplied by that public delivery
interface uses its original record even when its outer hints are
wrong. No corrected submission or acceptance-retry deadline is inferred.

A concrete encoder and decoder must preserve these logical values and
source denotations, including independently represented result objects
beside invalid hints or auxiliary evidence. The specification supplies
no blanket compatibility theorem for an unspecified legacy codec,
private-only resolver or arbitrary damaged bytes. Missing payload or
required ancestry was already excluded from complete results. Arbitrary
ciphertext without defined EXPORT semantics is not a logical tuple;
a documented old encoding of a tuple, however, requires a genuine
compatibility proof and cannot be dismissed as opaque.

The relevant-namespace visibility contract likewise is not bounded
discovery of arbitrary occurrences in unrelated namespaces. If the
phrase ``already recorded'' is given a broader implementation meaning
that discharges delivery for an occurrence outside the specified
public reader interface, an additional discovery argument is needed.
No such broader claim is used here. A wrong outer chain hint in a
record actually supplied through the existing interface is different:
the inventory still exposes its objects, whose signed domain determines
whether this interpreter may use them. Binary compatibility, external
access, native preparation and request inclusion are not consequences
of the structural lemma.


\paragraph{Public result kernels and DATA identity.}
A \defterm{public result kernel} contains a direct FINAL root for DATA
$D$, its exact DATA bytes, and the actual hash-linked DATA ancestry
needed to cover the proposed next DATA and connect it to the current
accepted anchor. It uses the fixed incoming boundary and voting snapshot.
\tagg{CheckFINAL} establishes intrinsic regional validity under
\cref{sec:receipt-kernel}; every independent global acceptance check below
still applies. A later session's receipt does not cover an earlier
session's DATA. The already accepted anchor keeps its original provenance.

The logical carrier may contain a larger complete regional result,
including full admission, recovery and native evidence. Those objects
remain available to regional validators and other interfaces, but are
not additional public-kernel dependencies. In particular, a publisher's
alternative PROPOSE or $\tagg{AUTH}_{r,g}$ path cannot replace the honest
signer's actual validation history or make the public interpreter fetch
that alternative. No kernel certifies an unbound execution-output,
state-snapshot or DA sidecar.

Kernels may share actual objects and use the common exact references of
\cref{sec:record-completeness}. The normalized DATA identity is
\[
 \operatorname{id}(B)=(\text{chain},\beta(B),r(B),\Hash(B)).
\]
Changing sender, receipt encoding or signer set does not change this
identity. A declared output interval cannot hide a complete next DATA
supplied as ancestry. An absent DATA object does not become present
because its hash appears in a signature or dependency list.

Completeness is checked for the required public kernel, not for all
optional evidence in its carrier. Malformed or irrelevant independently
represented logical objects cannot veto another kernel. A bad required DATA link,
wrong FINAL signature or missing required payload invalidates that
kernel. Try another complete kernel of the same DATA when one fails;
an invalid first wrapper never creates a permanent duplicate marker.
The original complete-result publication input remains unchanged, as
stated in \cref{sec:transport-scope}.

\paragraph{Canonical extraction: $\tagg{DeriveNextDATA}$.}
Let $p=(q,n,j)$ identify an Espresso block, namespace transaction and
record. Process these records in that order after the block's prelude.
At the current state, inspect its normalized public kernels for DATA at
$h(G)+1$. Order candidate proof choices by the tuple
\[
 (\beta(B),r(B),\Hash(B),\Hash(D),\operatorname{pos}(\pi)),
\]
where $D$ is the direct-receipt root providing coverage and
$\operatorname{pos}(\pi)$ is the canonical path of the actual object in the logical carrier.
Hashes use their fixed byte order. Equal normalized proof objects are
deduplicated. Ignore lower DATA heights except as evidence; never jump
to a higher height. Recompute the next-height candidates after a commit.
This is a finite deterministic scan of the admitted logical input.
Required ancestry and exact-reference objects obey the same declared
input and complete-message limits; optional evidence never creates a
mandatory external fetch.

The incoming boundary fixes the only session and region permitted to
supply that next DATA. Within that session, regional safety makes two
valid committed next DATA compatible and hence identical. The proof
ordering chooses evidence, not a fork. Neither a later session's
receipt nor a lower hash can override the fixed incoming boundary.

\paragraph{Acceptance checks.}
For a candidate $B$, derive its actual base from $G$ and require:
\begin{enumerate}[label=(\roman*),itemsep=2pt]
\item its signed domain is this chain's and its producing branch
version is current and open;
\item its incoming boundary, historical selection and authority assign
it to the current export duty; its authenticated DATA ancestry extends
the exact current $G$, including hashes, heights and original provenance;
\item its height is $h(G)+1$, its direct or same-session covering root
passes \tagg{CheckFINAL}, its actual supplied DATA hash path reaches
$G$, and its positions and quota agree with that fixed duty;
\item all required Espresso selection references strictly precede
the containing block $J_q$, with their historical snapshots verified;
\item if Force created its version in $J_{q_f}$, then $q>q_f$.
\end{enumerate}
\tagg{CheckFINAL} establishes the intrinsic admission, parent/reference
and native input-validity predicates through
\cref{prop:final-kernel-soundness}. The public algorithm does not rerun an
arbitrary publisher-supplied admission witness or native sidecar. This
is an explicit public-checker refinement; regional Vote, Ready and
recovery retain complete validation. Exact executable state and the
$C_{\rm pub}$ execution/DA interface remain separately verified and do
not precede the logical first-DATA commit here.

No check compares a sender's expected-base hint with $G$. Compatibility
is checked against the actual signed path instead. A record overlapping
$G$ may therefore supply its next compatible suffix. DRAIN can occur in
supporting ancestry, but is never appended as application DATA.

\paragraph{Commit one DATA: $\tagg{CommitNextDATA}$.}
On successful checks, atomically append $B$, increment its producing
session's accepted count and its region's accepted service count, and
remove that region from $\mathcal O$ if present. Count this DATA once.
A continuing duty keeps its window and receives the new window count;
an empty owed set is not refilled merely because another DATA arrived.
If $B=T_{\beta,r}$, cancel that duty's window and install the pending
REGULAR boundary described below in the same transaction. Durably expose
the updated $G$ and accounting before attempting a later DATA or
successor operation. The already saved pre-payload $\chi_E(q)$ and
$\mathcal O_E(q)$ are unchanged.

Acceptance is atomic per DATA and its associated accounting. There is
no all-or-nothing outer EXPORT transaction. An invalid later candidate
cannot roll back an accepted prefix. Batching is an implementation
optimization only if it preserves these committed prefixes and does
not postpone their exposure behind an unresolved successor operation.

\paragraph{Commit terminal DATA before resolving its successor.}
\label{sec:terminal-commit}
The terminal transaction stores the retained terminal anchor and its
commitment evidence, source context $(i_b,\beta,r)$, unchanged beta,
target slot $r+1$, signed selection reference and the regular selection
lower bound. It marks the boundary \defterm{pending resolution}; no
successor or DATA window is fabricated. Its owed ledger is the prior
ledger after the producing region's discharge, without a new round.
This transaction requires the completed session's incoming boundary
and authority, not $\tagg{AUTH}_{r+1,0}$.

Only the subsequent $\tagg{ResolveBoundary}$ operation obtains that
target authority, runs $\tagg{PrepareOwed}(r+1)$ and verifies the
prescribed historical REGULAR selection. A missing authority pauses
this operation after the terminal transaction has committed. A proven
absence of a candidate also leaves the boundary unassigned, but is
different from missing evidence. Neither case reverses accepted DATA,
restores the completed window, or opens a new window.
With evidence complete and no candidate yet, mark this resolution
attempt complete and continue source replay; a later prelude resumes
the prescribed historical search for the still-pending boundary.

Resolution resumes at its original canonical transition exactly once.
It preserves the historical selector and the public opening timestamp
in \cref{sec:ready}; local retrieval time never starts a new grace
period. An already fixed speculative handoff is checked against this
same history. Once resolution succeeds, the same record may supply a
successor DATA if its own references precede $J_q$. Thus normal
multi-session records remain possible. If the terminal publication
block is the first possible selection position, successor DATA still
cannot be accepted in that block.

Input authentication is incremental at this dependency boundary. It
authenticates the complete namespace and controls, then the authorities
needed for the prelude and current DATA transaction. It does not eagerly
demand every authority used by later operations in the record. In
particular, a preliminary whole-record check cannot wait for
$\tagg{AUTH}_{r+1,0}$ before committing an otherwise valid terminal DATA.

\paragraph{Durable subpositions and replay.}
\label{sec:subpositions}
The public cursor is
\[
 \zeta=(q,n,j,\nu,\mathit{stage}),\qquad
 \mathit{stage}\in\{\tagg{prelude},\tagg{scan},\tagg{resolve},\tagg{done}\}.
\]
Here $(q,n,j)$ is the block and record position, and $\nu$ counts DATA
transactions committed from that record. The prelude uses the sentinel
$n=j=\nu=0$. The stages mean prelude, next-DATA scan, pending boundary
resolution, and record complete, respectively. The boundary record
retains its original canonical resolution position across later scans.
At a DATA commit it also binds the normalized record commitment and
the resulting pending-boundary state. The next-height rule and this
cursor fix the continuation without dependence on private caches.
An implementation may retain a larger validation index as auxiliary
state. Duplicate source delivery resumes the cursor; it never reruns
a committed count or prelude.

An unresolved canonical boundary pauses subsequent replay at that
subposition. Earlier DATA commits remain exposed and final even though
the entire record or block is not yet complete. Once evidence is
available, resume the same transition with its original timestamp.
No out-of-order application of later controls is needed. Two observers
at different subpositions hold prefix-comparable DATA, just as they do
at different complete block heights.

If no candidate passes and no genuine canonical boundary wait remains,
finish the record and continue. An independently incomplete public kernel
does not suspend other complete kernels. A record skipped for a genuine
gap gains no future append permission: the interpreter never revisits
a completed past record to accept new DATA retroactively. Its objects
may remain historical evidence and may be referenced by a later record.
Closed-version DATA never becomes appendable through either path.
An incomplete namespace download, unlike an incomplete candidate,
prevents processing the block at all.

\paragraph{Scope of the transport interpretation.}
\label{sec:transport-scope}
The frozen publication input is still a \emph{complete regional result}:
DATA with its direct or covering FINAL receipt and the ancestry,
boundaries and source evidence needed to validate it. Correct senders
continue to offer that complete original value $R_0$, as distinguished
from the augmented regional bundle in \cref{sec:input-report-carriage}.
They offer it as soon as its original dependencies are available;
publication does not wait for an unrelated missing REPORT sidecar.
This revision does not promise
$T_{\rm pub}$ for a weaker kernel-only input, a hash, a request, a READY
message or a receipt with missing payload. The public checker uses a
sufficient subset of an already complete result's evidence; it does not
weaken the service's input obligation.

\Cref{prop:carrier-preservation} concerns the logical structured EXPORT
value and its actual evidence graph. It preserves a complete result
already carried by such a value, including one with incorrect outer
base, version or interval hints. No corrected wrapper or new publication
is needed. An encoding must faithfully expose that logical value and
its actual objects; no binary compatibility theorem for an unspecified
legacy codec is asserted. Opaque ciphertext, an undiscoverable namespace
or omitted required payload is not supplied by this normalization.

A stale signed path, closed version or missing predecessor still cannot
append. Sound FINAL signatures establish intrinsic regional validity,
not current external permission or a new completeness claim for absent
bytes. The frozen delivery guarantee remains transport only. A complete
terminal result does not require its successor authority, and a genuine
$\tagg{ResolveBoundary}$ wait follows its durable terminal commit.

The byte-retention rules of \cref{sec:source-retention} address later
access to exact old public objects. They do not add an archive-service
bound or infer a fresh deadline for every L1 object. The public kernel
avoids optional recovery-authority and native-sidecar fetches; required
controls, fixed boundaries and exact DATA/source bytes remain explicit
inputs. Broader transport interpretations that do not provide the
original logical complete result remain outside the progress proof.
